% c99-contract-test/de.tex — Vertragstest (deutsche Leitfassung)
%
% Kein Kursinhalt.  Dieses Kapitel benutzt jedes Konstrukt aus
% docs/authoring/CONTRACT.md mindestens einmal und ist die Vorlage für
% scripts/course/check_levels.py und für den HTML-Konverter.  Der Stoff ist
% absichtlich physikfrei: Rechnen mit ganzzahligen 2x2-Matrizen, jede
% Aussage wird im Text ausgerechnet.  Gebaut wird es nur mit
% --chapters c99 oder --with-test-chapter.

\chapter{Vertragstest: Zwei-mal-zwei-Matrizen}\label{c99}

Dieses Kapitel gehört nicht zum Kurs. Es ist eine Prüfvorlage: Es benutzt jede
Bauform, die der Autorenvertrag erlaubt, und rechnet dabei nur mit ganzen Zahlen.
Der eigentliche Kurs beginnt in Kapitel~\ref{c00}; dieser Verweis führt absichtlich
aus dem Kapitel hinaus.

\section{Zwei Matrizen und ihre Produkte}\label{c99.s1}

Wir betrachten zwei feste Matrizen mit ganzzahligen Einträgen,
\begin{equation}\label{c99.eq.matrizen}\claim{c99.cl.1}{author-stipulated}
  A = \begin{pmatrix} 1 & 2 \\ 0 & 1 \end{pmatrix},
  \qquad
  B = \begin{pmatrix} 1 & 0 \\ 3 & 1 \end{pmatrix}.
\end{equation}
Diese Festlegung ist die erste der Prüfvorlage \srcref{c99-fixture}{STIPULATIONS.md, S1}.
Ein Produkt zweier Matrizen entsteht nach der Regel „Zeile mal Spalte“. Für die
beiden möglichen Reihenfolgen ergibt das
\begin{align}
  AB &= \begin{pmatrix} 1\cdot 1 + 2\cdot 3 & 1\cdot 0 + 2\cdot 1 \\
                        0\cdot 1 + 1\cdot 3 & 0\cdot 0 + 1\cdot 1 \end{pmatrix}
      = \begin{pmatrix} 7 & 2 \\ 3 & 1 \end{pmatrix},
  \label{c99.eq.ab}\claim{c99.cl.2}{derived} \\
  BA &= \begin{pmatrix} 1\cdot 1 + 0\cdot 0 & 1\cdot 2 + 0\cdot 1 \\
                        3\cdot 1 + 1\cdot 0 & 3\cdot 2 + 1\cdot 1 \end{pmatrix}
      = \begin{pmatrix} 1 & 2 \\ 3 & 7 \end{pmatrix}.
  \label{c99.eq.ba}\claim{c99.cl.3}{derived}
\end{align}

\begin{proposition}[Die Reihenfolge zählt]\label{c99.prop.reihenfolge}\claim{c99.cl.4}{derived}
  Für die Matrizen aus \eqref{c99.eq.matrizen} gilt $AB \neq BA$.
\end{proposition}
\begin{proof}
  Der Eintrag oben links ist in \eqref{c99.eq.ab} gleich $7$ und in
  \eqref{c99.eq.ba} gleich $1$.
\end{proof}

\begin{tpcfigure}{c99.fig.scherung}
  \demo{c99.demo.scherung}
  \caption{Die Matrix $A$ als Abbildung der Ebene: Der Punkt $(x,y)$ geht auf
    $(x+2y,\,y)$. Das gestrichelte Einheitsquadrat wird dabei zum schraffierten
    Parallelogramm mit den Ecken $(0,0)$, $(1,0)$, $(3,1)$ und $(2,1)$; die
    Grundseite bleibt liegen, die Deckseite wandert um zwei Einheiten nach rechts.
    Diese Bildunterschrift ist absichtlich lang, damit der Zeilenumbruch langer
    deutscher Zusammensetzungen wie Parallelogrammflächeninhaltsberechnung und
    Einheitsquadratverschiebungsabbildung in Druck und Netzfassung geprüft wird.}
\end{tpcfigure}

\tpcdef{T1}

Mit \defref{T1} liest man aus \eqref{c99.eq.ab} und \eqref{c99.eq.ba} ab:
\begin{equation}\label{c99.eq.spuren}\claim{c99.cl.5}{derived}
  \tr(AB) = 7 + 1 = 8,
  \qquad
  \tr(BA) = 1 + 7 = 8.
\end{equation}
Der Mittelwert der vier Einträge von $AB$ ist \num[c99.cl.19]{3{,}25}{}.
Die Produkte sind verschieden, ihre Spuren aber gleich. Das ist kein Zufall dieser
beiden Matrizen\zoomref{c99.v1}:

\begin{theorem}[Spur eines Produkts]\label{c99.thm.spur}\claim[proof-linked]{c99.cl.6}{derived}
  Für alle $2\times 2$-Matrizen $X$ und $Y$ gilt $\tr(XY) = \tr(YX)$.
\end{theorem}

\begin{vertiefung}[id=c99.v1,
    titel={Warum die Spuren übereinstimmen},
    wozu={Der Kern stellt Theorem~\ref{c99.thm.spur} ohne Beweis hin; hier wird es
      für beliebige Einträge ausgerechnet.},
    voraussetzung={Kern~\ref{c99.s1} bis Gleichung~\eqref{c99.eq.spuren}; nicht
      nötig für Abschnitt~\ref{c99.s2}.},
    zurueck={Die Spur eines Produkts hängt nicht von der Reihenfolge der Faktoren
      ab, das Produkt selbst schon.}]
  Seien $X$ und $Y$ Matrizen mit Einträgen $x_{jk}$ und $y_{jk}$. Die
  Diagonaleinträge der beiden Produkte sind
  \begin{align}
    (XY)_{11} &= x_{11}y_{11} + x_{12}y_{21}, &
    (XY)_{22} &= x_{21}y_{12} + x_{22}y_{22},
    \label{c99.eq.diag-xy}\claim{c99.cl.7}{derived} \\
    (YX)_{11} &= y_{11}x_{11} + y_{12}x_{21}, &
    (YX)_{22} &= y_{21}x_{12} + y_{22}x_{22}.
    \label{c99.eq.diag-yx}\claim{c99.cl.8}{derived}
  \end{align}

  \begin{lemma}[Spur eines Produkts, ausgerechnet]\label{c99.lem.spur}\claim{c99.cl.9}{derived}
    Für alle $2\times 2$-Matrizen $X$ und $Y$ gilt
    \begin{equation}\label{c99.eq.spur-allgemein}
      \tr(XY) = x_{11}y_{11} + x_{12}y_{21} + x_{21}y_{12} + x_{22}y_{22} = \tr(YX).
    \end{equation}
  \end{lemma}
  \begin{proof}
    Die Summe der beiden Einträge in \eqref{c99.eq.diag-xy} ist der mittlere
    Ausdruck. Die Summe der beiden Einträge in \eqref{c99.eq.diag-yx} besteht aus
    denselben vier Produkten in anderer Reihenfolge.
  \end{proof}

  \tpcdef{T2}

  Für die beiden Matrizen des Kerns ist nach \eqref{c99.eq.ab} und \eqref{c99.eq.ba}
  \begin{equation}\label{c99.eq.kommutator}\claim{c99.cl.10}{derived}
    \comm{A}{B} = \begin{pmatrix} 7-1 & 2-2 \\ 3-3 & 1-7 \end{pmatrix}
                = \begin{pmatrix} 6 & 0 \\ 0 & -6 \end{pmatrix},
    \qquad
    \tr\comm{A}{B} = 6 - 6 = 0.
  \end{equation}

  \begin{tpcfigure}{c99.fig.kommutator}
    \caption{Der Kommutator aus \eqref{c99.eq.kommutator} als Abbildung: Er streckt
      um den Faktor $6$ und spiegelt an der waagerechten Achse.}
  \end{tpcfigure}

  Dass die Spur Summen respektiert, wird im Fundament nachgerechnet\zoomref{c99.f2}.
\end{vertiefung}

Der nächste Absatz ist für Leserinnen und Leser geschrieben, die den Block
übersprungen haben: Er benutzt nur \eqref{c99.eq.spuren} und Theorem~\ref{c99.thm.spur}.

\begin{tpcfigure}{c99.fig.scherung-b}
  \demo{c99.demo.scherung-b}
  \caption{Die Matrix $B$ als Abbildung: $(x,y)$ geht auf $(x,\,3x+y)$. Zweites
    Bild mit Demo auf derselben Seite.}
\end{tpcfigure}

\begin{tpctable}{c99.tab.produkte}
  \caption{Die Einträge der beiden Produkte aus \eqref{c99.eq.ab} und \eqref{c99.eq.ba}.}
  \begin{tabular}{lcccc}
    \toprule
    Produkt & $(1,1)$ & $(1,2)$ & $(2,1)$ & $(2,2)$ \\
    \midrule
    $AB$ & $7$ & $2$ & $3$ & $1$ \\
    $BA$ & $1$ & $2$ & $3$ & $7$ \\
    \bottomrule
  \end{tabular}
\end{tpctable}

\begin{merksatz}[Merksatz]
  Bei Matrizen kommt es auf die Reihenfolge der Faktoren an.
\end{merksatz}

\section{Scherungen}\label{c99.s2}

\subsection{Eine Familie von Matrizen}\label{c99.s2.1}

Für jede ganze Zahl $t$ sei
\begin{equation}\label{c99.eq.scherung}\claim{c99.cl.11}{author-stipulated}
  M(t) = \begin{pmatrix} 1 & t \\ 0 & 1 \end{pmatrix}.
\end{equation}
Die Matrix $A$ aus \eqref{c99.eq.matrizen} ist $M(2)$.

\begin{theorem}[Scherungen addieren sich]\label{c99.thm.scherung}\claim[proof-included]{c99.cl.12}{derived}
  Für alle ganzen Zahlen $s$ und $t$ gilt $M(s)M(t) = M(s+t)$. Für jede ganze Zahl
  $n \geq 1$ gilt außerdem $M(t)^n = M(nt)$.
\end{theorem}
\begin{proof}
  Die erste Aussage ist eine Multiplikation:
  \begin{equation}\label{c99.eq.scherung-produkt}
    M(s)M(t)
    = \begin{pmatrix} 1\cdot 1 + s\cdot 0 & 1\cdot t + s\cdot 1 \\
                      0\cdot 1 + 1\cdot 0 & 0\cdot t + 1\cdot 1 \end{pmatrix}
    = \begin{pmatrix} 1 & s+t \\ 0 & 1 \end{pmatrix}
    = M(s+t).
  \end{equation}
  Die zweite folgt durch vollständige Induktion über $n$. Für $n = 1$ ist nichts
  zu zeigen. Der Induktionsschritt ist ein eigenes Lemma.
  \begin{lemma}[Induktionsschritt]\label{c99.lem.schritt}\claim{c99.cl.13}{derived}
    Gilt $M(t)^n = M(nt)$, so gilt auch $M(t)^{n+1} = M((n+1)t)$.
  \end{lemma}
  \begin{proof}
    Nach Voraussetzung ist $M(t)^{n+1} = M(t)^n M(t) = M(nt)M(t)$, und nach
    \eqref{c99.eq.scherung-produkt} mit $s = nt$ ist das $M(nt + t) = M((n+1)t)$.
  \end{proof}
  Mit Lemma~\ref{c99.lem.schritt} gilt die Aussage für alle $n \geq 1$.
\end{proof}

\begin{corollary}[Scherungen vertauschen]\label{c99.cor.vertauschen}\claim{c99.cl.18}{derived}
  Für alle ganzen Zahlen $s$ und $t$ gilt $M(s)M(t) = M(t)M(s)$.
\end{corollary}
\begin{proof}
  Nach Theorem~\ref{c99.thm.scherung} sind beide Seiten gleich $M(s+t)$, denn
  $s + t = t + s$.
\end{proof}

\begin{example}[Zahlenbeispiel]\label{c99.ex.zahlen}
  Mit $s = 2$ und $t = 3$ liefert \eqref{c99.eq.scherung-produkt} die Gleichung
  $M(2)M(3) = M(5)$, und mit $t = 2$, $n = 3$ ergibt sich $A^3 = M(2)^3 = M(6)$.
\end{example}

\begin{remark}
  Die Menge aller $M(t)$ ist unter Multiplikation abgeschlossen. Was daraus
  folgt, steht im Fundament\zoomref{c99.f1}.
\end{remark}

\subsection{Eine Fallunterscheidung}\label{c99.s2.2}

Der Betrag des Eintrags oben rechts von $M(t)$ ist
\begin{equation}\label{c99.eq.betrag}\claim{c99.cl.14}{author-stipulated}
  f(t) =
  \begin{cases}
    t  & \text{\tpcifenglish{if}{falls}}\ t \geq 0, \\
    -t & \text{\tpcifenglish{if}{falls}}\ t < 0.
  \end{cases}
\end{equation}

\begin{proposition}[Der Betrag ist nicht negativ]\label{c99.prop.betrag}\claim{c99.cl.15}{derived}
  Für jede ganze Zahl $t$ gilt $f(t) \geq 0$ und $f(t)^2 = t^2$.
\end{proposition}
\begin{proof}
  Wir unterscheiden die beiden Fälle aus \eqref{c99.eq.betrag}.
  \begin{description}
    \item[Fall 1: $t \geq 0$.] Dann ist $f(t) = t \geq 0$ und $f(t)^2 = t^2$.
    \item[Fall 2: $t < 0$.] Dann ist $f(t) = -t > 0$ und
      $f(t)^2 = (-t)^2 = t^2$.
  \end{description}
\end{proof}

\begin{intuition}
  Eine Scherung verschiebt jede waagerechte Linie in sich; nur der Betrag der
  Verschiebung wächst mit der Höhe. Die Kästen für Elementtypen sind getönt, die
  Tiefenblöcke nicht.
\end{intuition}

\begin{exercise}[Die Umkehrung]\label{c99.exc.umkehrung}
  Berechne $M(t)M(-t)$.
  \begin{enumerate}
    \item Benutze \eqref{c99.eq.scherung-produkt}.
    \item Vergleiche mit der direkten Multiplikation.
  \end{enumerate}
\end{exercise}
\begin{solution}
  Mit $s = t$ und $-t$ an der Stelle von $t$ liefert \eqref{c99.eq.scherung-produkt}
  die Matrix $M(t - t) = M(0)$, und $M(0)$ hat die Einträge $1, 0, 0, 1$.
  \begin{itemize}
    \item Direkt: Der Eintrag oben rechts ist $1\cdot(-t) + t\cdot 1 = 0$.
    \item Die übrigen Einträge sind wie in \eqref{c99.eq.scherung-produkt}.
  \end{itemize}
\end{solution}

\begin{fundament}[id=c99.f1,
    titel={Die Scherungen bilden eine Gruppe},
    wozu={Der Kern benutzt nur das Produkt zweier Scherungen; hier wird
      zusammengetragen, was daraus für die ganze Familie folgt.},
    voraussetzung={Kern~\ref{c99.s2.1}, Theorem~\ref{c99.thm.scherung}.},
    zurueck={Jede Scherung $M(t)$ lässt sich durch $M(-t)$ rückgängig machen.}]
  \begin{proposition}[Einselement und Inverses]\label{c99.prop.gruppe}\claim{c99.cl.16}{derived}
    Für alle ganzen Zahlen $t$ gilt $M(0)M(t) = M(t)$ und $M(t)M(-t) = M(0)$,
    ebenso mit vertauschten Faktoren.
  \end{proposition}
  \begin{proof}
    Beides ist Theorem~\ref{c99.thm.scherung} mit einer besonderen Wahl von $s$:
    \begin{equation}\label{c99.eq.gruppe}
      M(0)M(t) = M(0 + t) = M(t),
      \qquad
      M(t)M(-t) = M(t - t) = M(0).
    \end{equation}
    Die Gleichungen mit vertauschten Faktoren folgen genauso, weil $s + t = t + s$.
  \end{proof}
\end{fundament}

\begin{fundament}[id=c99.f2,
    titel={Die Spur respektiert Summen},
    wozu={Die Vertiefung benutzt, dass die Spur einer Differenz die Differenz der
      Spuren ist; hier steht die Rechnung.},
    voraussetzung={\defref{T1}.},
    zurueck={Aus $\tr(AB) = \tr(BA)$ folgt deshalb $\tr(AB - BA) = 0$.}]
  Für Matrizen $X$ und $Y$ mit Einträgen $x_{jk}$ und $y_{jk}$ und eine ganze Zahl
  $c$ ist
  \begin{equation}\label{c99.eq.spur-linear}\claim{c99.cl.17}{derived}
    \tr(X + cY) = (x_{11} + c\,y_{11}) + (x_{22} + c\,y_{22})
                = \tr X + c\,\tr Y.
  \end{equation}
  Mit $c = -1$ ist das die Aussage über Differenzen.
\end{fundament}
